% !TEX program = xelatex
% [EN] The English note uses the same signed bends and beam-frame convention as the Chinese version. / [CN] 英文稿与中文稿采用相同的带符号弯角及局部束流坐标约定。
\documentclass[11pt,a4paper]{article}
\usepackage[margin=23mm,headheight=15pt]{geometry}
\usepackage{fontspec}
\setmainfont{TeX Gyre Termes}
\setsansfont{TeX Gyre Heros}
\usepackage{amsmath,amssymb,bm,booktabs,tabularx,array,xcolor,fancyhdr,titlesec}
\usepackage{xurl}
\usepackage[colorlinks=true,linkcolor=blue!45!black,citecolor=blue!45!black,urlcolor=blue!45!black]{hyperref}
\definecolor{ink}{HTML}{193B55}
\titleformat{\section}{\large\bfseries\sffamily\color{ink}}{\thesection}{0.7em}{}
\titleformat{\subsection}{\normalsize\bfseries\sffamily\color{ink}}{\thesubsection}{0.7em}{}
\linespread{1.03}
\setlength{\parskip}{3pt}
\setlength{\emergencystretch}{2em}
\renewcommand{\arraystretch}{1.18}
\newcolumntype{Y}{>{\raggedright\arraybackslash}X}
\newcommand{\pvec}{\bm p}
\newcommand{\Pt}{\mathsf P}
\pagestyle{fancy}
\fancyhf{}
\fancyhead[L]{\small\sffamily A polarimeter downstream of SRC}
\fancyhead[R]{\small BigDpol/T11 and F13}
\fancyfoot[C]{\thepage}
\hypersetup{pdftitle={Why Install a Polarimeter Downstream of SRC?},pdfauthor={DPOLAR Project}}
\begin{document}
\thispagestyle{empty}
\begin{center}
{\LARGE\sffamily\color{ink}Why Install a Polarimeter Downstream of SRC?}\par\vspace{10pt}
DPOLAR Project\quad---\quad Sources checked: 8 September 2026
\end{center}
\vspace{5pt}
\noindent
A polarimeter at the former BigDpol/T11 location could use higher local beam intensity to give faster statistical feedback during polarization tuning. Comparison with F13 would help localize polarization changes along the beam line. The known spin rotation between the stations also provides complementary measurement projections, enabling more complete polarization reconstruction.

\section{Higher upstream rates give faster tuning feedback}

The most practical reason for installing a polarimeter after SRC, before BigRIPS, is the potentially higher beam intensity available there. The official SRC table lists a planning current of $30\,\mathrm{pnA}$ for polarized deuterons at $250\,\mathrm{MeV/u}$, subject to source development; the published total-rate limit at F3 is $10^7\,\mathrm{s^{-1}}$~\cite{current,limit}. These are different constraints---accelerator delivery and downstream total rate---and the planning current does not establish the available intensity at $190\,\mathrm{MeV/u}$.

For fixed target thickness, acceptance, analyzing powers and live-time fraction,
\begin{equation}
 N\propto I_{\mathrm b}t,\qquad
 \sigma(P)\propto N^{-1/2},\qquad
 t_{\sigma(P)=\mathrm{const}}\propto I_{\mathrm b}^{-1}.
\end{equation}
Within the detector's linear counting range, higher local intensity therefore shortens the time needed to tune PIS/RF modes, Wien-filter orientation and accelerator extraction or beam conditions at a given statistical precision.

The allowed high-intensity operating mode at T11 must be confirmed with the accelerator and BigRIPS teams, including any local beam termination or attenuation. The 2025 BIS report describes primary-beam permission based on attenuators upstream of the accelerator complex~\cite{bis}; it does not establish an approved high-T11/low-F13 attenuation scheme.

\section{Two stations localize polarization changes}

T11 provides a reference after the source, accelerators, SRC and upstream transport. F13 measures the incident beam after DMT2/DMT3 and BigRIPS, close to the experiment and upstream of the SAMURAI reaction target:
\begin{equation*}
\begin{aligned}
\mathrm{SRC}\to DMT1\to\boxed{\mathrm{BigDpol/T11}}
&\to DMT2,DMT3\to F0\\[-1pt]
&\to\mathrm{BigRIPS}\ (F0\text{--}F7)\to\boxed{F13}.
\end{aligned}
\end{equation*}
The SAMURAI analyzing magnet is downstream of the target and is excluded from this incident-beam transport~\cite{local,geometry}.

For the same source mode and comparable beam conditions, measure at T11, apply the known transport model and compare the predicted observables with F13. Because one station has blind components, this prediction requires sufficient source-state constraints; unmeasured components must not be set to zero by default. A general density matrix requires a joint fit to both stations.

\noindent\begin{minipage}{\linewidth}
\begingroup\small\noindent
\begin{tabularx}{\linewidth}{@{}p{0.36\linewidth}Y@{}}
\toprule
Observation & Where to investigate first \\
\midrule
T11 differs from its reference & Source mode, spin orientation, extraction and transport upstream of T11.\\
T11 agrees with its reference; F13 differs from the prediction & Transport fields, trajectories and beam selection between stations; F13 instrument response.\\
Both change consistently with the transport model & A common upstream change.\\
\bottomrule
\end{tabularx}\par
\endgroup
\vspace{4pt}

Account for efficiency, live time, beam position and acceptance. Slits, attenuation and other phase-space selection can change the sampled beam and must be included in the comparison. Two stations help localize the problematic section; they do not uniquely identify a faulty magnet.
\end{minipage}\par\vspace{1ex}

\section{Two stations provide complementary spin projections}\label{sec:projections}

An additional benefit is sensitivity to polarization components that one station cannot measure. At $190\,\mathrm{MeV/u}$, the current central-trajectory model gives a relative spin rotation of about $27.5^\circ$ between the local beam frames at T11 and F13. The two stations therefore measure complementary projections~\cite{local,geometry,bmt,constants}.

In the present multi-angle L/R/U/D geometry, one station has blind vector and tensor components. The known, non-degenerate relative rotation removes these blind directions in the combined response: in principle, the two stations can identify all eight real parameters of a normalized spin-1 density matrix, or all five tensor components in a tensor-only analysis~\cite{local}. This requires calibrated responses, non-degenerate analyzing-power combinations across polar angles, constrained normalization and efficiencies, and a known transport model relating the sampled beam states. Full rank establishes identifiability, not a guarantee of statistical precision. Appendix~\ref{app:transport} gives the compact derivation.

\section{Summary}

The first reason to install a polarimeter at T11 is faster polarization-tuning feedback if higher local intensity is permitted. Its comparison with F13 would then help distinguish upstream changes from changes between the stations. The known spin rotation adds a third benefit: complementary projections for complete polarization reconstruction under the stated measurement conditions.

\clearpage
\appendix
\numberwithin{equation}{section}
\section{Spin transport and measurement rank}\label{app:transport}

\subsection{Central-trajectory rotation at 190 MeV/u}

Take local $z$ along the beam, $y$ vertical and $x$ completing a right-handed frame. DMT2/DMT3 each bend by $-50^\circ$; BigRIPS D1--D6 contribute $(-30,-30,-30,+30,+30,-30)^\circ$; F7 to the pre-target F13 station adds no net bend~\cite{local,geometry}. Thus the signed net orbital bend, rather than the sum of absolute bend angles, is
\begin{equation}
 \Theta_{\mathrm{T11}\to F13}=-50^\circ-50^\circ-60^\circ+0^\circ=-160^\circ.
\end{equation}
For approximately constant energy and a horizontal central trajectory dominated by transverse dipole fields, Thomas--BMT gives~\cite{bmt,constants}
\begin{align}
 \psi_{\mathrm{lab}}&=(1+G_d\gamma)\Theta,&
 \delta&=\psi_{\mathrm{lab}}-\Theta=G_d\gamma\Theta,\\
 G_d&=-0.14298726968,&
 \gamma&=1+\frac{2T_u}{m_dc^2},\quad m_dc^2=1875.612945\,\mathrm{MeV}.
\end{align}
For kinetic energy per nucleon $T_u=190\,\mathrm{MeV/u}$, $\gamma=1.20260044$, giving
\begin{equation}
 \psi_{\mathrm{lab}}=-132.487^\circ,\qquad \boxed{\delta=+27.513^\circ.}
\end{equation}
The relevant angle is $\delta$ between local beam frames; energy changes or significant energy loss require recalculation.

\subsection{Shared rotation and removal of blind components}

With $c=\cos\delta$ and $s=\sin\delta$, the vector and symmetric traceless tensor obey
\begin{equation}
 R_y=\begin{pmatrix}c&0&s\\0&1&0\\-s&0&c\end{pmatrix},\qquad
 \pvec_{13}=R_y\pvec_{\mathrm{T11}},\qquad
 \Pt_{13}=R_y\Pt_{\mathrm{T11}}R_y^{\mathsf T}.
\end{equation}
Vector directions and tensor principal axes share this spatial rotation: an initially longitudinal axis becomes $(s,0,c)=(0.461951,0,0.886906)$, while a vertical axis is unchanged. A tensor axis is defined modulo $180^\circ$. Double-angle terms, such as
\begin{equation}
 p'_{xz}=\tfrac12\sin2\delta\,(p_{zz}-p_{xx})+\cos2\delta\,p_{xz},
\end{equation}
describe tensor components, not a doubled axis rotation; an arbitrary axis transforms as $R_y\hat{\bm n}$.

A spin-1 density matrix has three vector and five tensor parameters. At one multi-angle L/R/U/D station, $p_z$ is absent from the scattering response and the $\sin2\phi$ term containing $p_{xy}$ vanishes at the cardinal azimuths~\cite{local}. At the second station,
\begin{equation}
 p'_x=c\,p_x+s\,p_z,\qquad p'_{yz}=-s\,p_{xy}+c\,p_{yz}.
\end{equation}
Since $s=0.461951\ne0$, both blind components enter observable channels. Under the calibration and transport conditions stated in Section~\ref{sec:projections},
\begin{center}
\begin{tabular}{@{}lcc@{}}
\toprule
Multi-angle response & Vector + tensor & Tensor only \\
\midrule
One station, or two without relative rotation & $6/8$ & $4/5$ \\
T11 and F13 with known $R_y(27.513^\circ)$ & $8/8$ & $5/5$ \\
\bottomrule
\end{tabular}
\end{center}
One station can already measure a tensor amplitude along a known vertical axis. General reconstruction needs independent projections: for the present identical azimuthal layouts, the intervening spin rotation supplies them; a controlled rotation or another suitable geometry could also do so.

\begingroup\footnotesize
\setlength{\parskip}{0pt}
\begin{thebibliography}{8}
\setlength{\itemsep}{3pt}
\bibitem{local} DPOLAR, \emph{Deuteron Spin Density Matrix, Beam Transport, and Polarimeter Identifiability}. Detailed station locations, conventions and response derivations: \path{docs/physics/tensor_identifiability/deuteron_spin_density_matrix/main_en.tex}.
\bibitem{geometry} K. Sekiguchi et al., \emph{Plan of Measurement of Polarized Deuteron--Proton Scattering at RIBF}, RIKEN APR 42 (2009), 27; K. Kusaka et al., \emph{Present status of the beam transport line from SRC to BigRIPS}, RIKEN APR 51 (2018), 173, \url{https://www.nishina.riken.jp/researcher/APR/APR051/pdf/173.pdf}.
BigRIPS and SAMURAI geometry: \url{https://ribf.riken.jp/BigRIPSInfo/}; \href{https://ribf.riken.jp/BigRIPSInfo/lise/NoUCSF_Be_fragmentation_Standard_SAMURAI.lpp}{NoUCSF\_Be\_fragmentation\_Standard\_SAMURAI.lpp}, used only for geometry, not as a polarized-deuteron run setting.
\bibitem{current} RIKEN, \emph{Accelerator Technical Information}, SRC table and source-development qualification: \url{https://www.nishina.riken.jp/ribf/accelerator/tecinfo.html}.
\bibitem{limit} RIKEN, \emph{Secondary beam intensity expected}, total rate at F3: \url{https://www.nishina.riken.jp/ribf/BigRIPS/intensity.html}.
\bibitem{bis} M. Komiyama et al., \emph{Development of a machine protection beam interlock system for the RIBF}, RIKEN APR 58 (2025), 98--99: \url{https://www.nishina.riken.jp/researcher/APR/APR058/pdf/98.pdf}.
\bibitem{bmt} V. Bargmann, L. Michel and V. L. Telegdi, Phys. Rev. Lett. 2 (1959), 435--436, \href{https://doi.org/10.1103/PhysRevLett.2.435}{doi:10.1103/PhysRevLett.2.435}.
\bibitem{constants} P. J. Mohr et al., \emph{CODATA Recommended Values of the Fundamental Physical Constants: 2022}, Rev. Mod. Phys. 97 (2025), 025002, \href{https://doi.org/10.1103/RevModPhys.97.025002}{doi:10.1103/RevModPhys.97.025002}. Numerical implementation: \path{analysis/spin_transport/bmt.py}.
\end{thebibliography}
\endgroup
\end{document}
